(1) 约定 \(\mathsf J\mathbf v=\hat{\mathbf z}\times\mathbf v\)，因此对沿 \(\pm z\) 的平面波，
\(\mathbf H_\pm=\pm Z_0^{-1}\mathsf J\mathbf E_\pm\)。
在薄层处令 \(\mathbf E_r,\mathbf E_t\) 为反射和透射复振幅。切向电场连续给出
\(\mathbf E_t=\mathbf E_i+\mathbf E_r\)。
磁场跳跃条件逐项写成
\[
\mathsf J\left[\frac{\mathsf J\mathbf E_t}{Z_0}
-\frac{\mathsf J(\mathbf E_i-\mathbf E_r)}{Z_0}\right]
=\boldsymbol{\sigma}\mathbf E_t.
\]
利用 \(\mathsf J^2=-\mathbb I\) 得到
\[
-\frac{2}{Z_0}\mathbf E_r=\boldsymbol{\sigma}\mathbf E_t,\qquad
(\mathbb I+\tfrac{Z_0}{2}\boldsymbol{\sigma})\mathbf E_t=\mathbf E_i.
\]
记 \(D=(1+a)^2+b^2>0\)，故
\[
\boxed{\mathcal T=[(1+a)\mathbb I+b\mathsf J]^{-1}
=\frac{(1+a)\mathbb I-b\mathsf J}{D}},\qquad
\boxed{\mathcal R=\mathcal T-\mathbb I=-[a\mathbb I+b\mathsf J]\mathcal T}.
\]
对实验室固定的圆偏振列矢 \(\mathbf e_s=(1,-is)^T/\sqrt2\)（\(s=\pm1\)），
有 \(\mathsf J\mathbf e_s=is\mathbf e_s\)。所以
\[
\boxed{t_s=\frac1{1+a+is b}},\qquad
\boxed{r_s=-\frac{a+is b}{1+a+is b}}.
\]
圆偏振的 \(s\) 标签以 \(\mathsf J\) 的本征值定义；若改用相反的圆偏振标签或 \(e^{+i\omega t}\) 约定，相位符号必须同时更改，不可直接套用旋转角符号。

(2) 对纯霍尔层 \(a=0\)，
\[
\boxed{\mathbf E_t=\frac{E_0}{1+b^2}(\hat{\mathbf x}-b\hat{\mathbf y})},\qquad
\boxed{\mathbf E_r=-\frac{E_0}{1+b^2}(b^2\hat{\mathbf x}+b\hat{\mathbf y})}.
\]
因所有分量在所选相位约定下均为实数，出射波仍为线偏振。
透射偏振角按 \(x\) 朝 \(y\) 为正，取
\[
\boxed{\theta_F=-\arctan b\quad(\bmod\ \pi)}.
\]
对 \(b\ne0\)，直接计算
\(\mathbf E_t\cdot\mathbf E_r
=E_0^2(-b^2+b^2)/(1+b^2)^2=0\)，
故反射线偏振方向为 \(\theta_F+\pi/2\pmod\pi\)。
时间平均 Poynting 通量 \(S=|\mathbf E|^2/(2Z_0)\)，于是
\[
\boxed{T=\frac1{1+b^2}},\qquad
\boxed{R=\frac{b^2}{1+b^2}},\qquad R+T=1.
\]
薄层的霍尔电流 \(\mathbf K=\sigma_H\mathsf J\mathbf E_t\)
与层内电场垂直，故 \(\langle\mathbf K\cdot\mathbf E_t\rangle=0\)，
无耗散。\(b\to0\) 时 \(T\to1,R\to0,\theta_F\to0\)；
恰在 \(b=0\) 时反射振幅为零，因而不存在可定义的反射偏振方向。

(3) 由(1)式，对线偏振入射
\[
\mathbf E_t=\frac{E_0}{D}[(1+a)\hat{\mathbf x}-b\hat{\mathbf y}],\qquad
\mathbf E_r=-\frac{E_0}{D}[(a+a^2+b^2)\hat{\mathbf x}+b\hat{\mathbf y}].
\]
直接平方并化简，得到
\[
\boxed{T=\frac1D},\qquad
\boxed{R=\frac{a^2+b^2}{D}}.
\]
平均面焦耳热为
\(\langle p_{\rm J}\rangle
=\tfrac12\operatorname{Re}(\mathbf E_t^*\cdot\boldsymbol{\sigma}\mathbf E_t)
=\tfrac12\sigma_d|\mathbf E_t|^2\)，
因为反对称实矩阵 \(\mathsf J\) 的二次型不贡献实功。
除以入射能流 \(E_0^2/(2Z_0)\)，
\[
\boxed{A=\frac{Z_0\sigma_d}{D}=\frac{2a}{(1+a)^2+b^2}\ge0}.
\]
于是 \(T+R+A=[1+a^2+b^2+2a]/D=1\)。
固定 \(b\)，
\[
\frac{\partial A}{\partial a}
=\frac{2(1+b^2-a^2)}{[(1+a)^2+b^2]^2}.
\]
在 \(a\ge0\) 上导数只于
\(a_*=\sqrt{1+b^2}\) 由正转负，且 \(A(0)=0\)、\(A(a\to\infty)=0\)，
因此唯一全局最大值为
\[
\boxed{\sigma_{d,*}=\frac{2\sqrt{1+b^2}}{Z_0}},\qquad
\boxed{A_{\max}=\frac1{1+\sqrt{1+b^2}}\le\frac12}.
\]
无霍尔效应 \(b=0\) 时恢复自由空间单张电阻片的 \(50\%\) 吸收上界；
\(a\to\infty\) 时 \(R\to1,T,A\to0\)，体现理想反射极限。
\(\sigma_d,\sigma_H\) 的量纲均为西门子（二维片电导），
\(Z_0\sigma\)、\(a,b,D,T,R,A\) 均无量纲。

(4) 对无耗非磁性介质，正、反行波的磁场分别为
\(\pm(n_j/Z_0)\mathsf J\mathbf E\)。跳跃条件与
\(\mathbf E_t=\mathbf E_i+\mathbf E_r\) 联立后得到
\[
2n_1\mathbf E_i=[(n_1+n_2)\mathbb I+Z_0\boldsymbol{\sigma}]\mathbf E_t.
\]
令 \(q=n_1+n_2>0\)、\(h=Z_0\sigma_H\)，纯霍尔层的矩阵为
\[
\boxed{\mathcal T_{12}=\frac{2n_1(q\mathbb I-h\mathsf J)}{q^2+h^2}},\qquad
\boxed{\mathcal R_{12}=\mathcal T_{12}-\mathbb I}.
\]
对 \(x\) 线偏振入射，透射偏振角
\[
\boxed{\theta_{F,12}=-\arctan\frac{h}{n_1+n_2}}.
\]
此时应按不同介质的波阻抗计算能流，而非直接用电场振幅平方作透过率：
\[
\boxed{T_{12}=\frac{n_2}{n_1}
\frac{|\mathbf E_t|^2}{|\mathbf E_i|^2}
=\frac{4n_1n_2}{(n_1+n_2)^2+h^2}},\qquad
\boxed{R_{12}=\frac{(n_1-n_2)^2+h^2}{(n_1+n_2)^2+h^2}}.
\]
二者相加恰为1，符合纯霍尔层无焦耳耗散。
由 \(Z_0=1/(\epsilon_0c)\) 和 \(h_{\rm P}=2\pi\hbar\)，
\[
Z_0\frac{e^2}{h_{\rm P}}
=\frac{e^2}{2\pi\epsilon_0\hbar c}=2\alpha_{\rm fs}.
\]
故整数霍尔平台的旋转角为
\[
\boxed{\theta_{F,12}
=-\arctan\!\left(\frac{2\nu\alpha_{\rm fs}}{n_1+n_2}\right)}.
\]
当 \(\sigma_H=0\) 时恢复通常的
\(T_{12}=4n_1n_2/(n_1+n_2)^2\)、
\(R_{12}=(n_1-n_2)^2/(n_1+n_2)^2\)；
当 \(n_1=n_2=1\) 时
\(\theta_F=-\arctan(\nu\alpha_{\rm fs})
=-\nu\alpha_{\rm fs}+O((\nu\alpha_{\rm fs})^3)\)。
这里的量子化关系假定低频、无耗、局域的霍尔响应，
不可无条件外推到有限频率吸收共振或有损基底。